% Options for packages loaded elsewhere
\PassOptionsToPackage{unicode}{hyperref}
\PassOptionsToPackage{hyphens}{url}
\PassOptionsToPackage{dvipsnames,svgnames,x11names}{xcolor}
%
\documentclass[
  10pt,
  a4paper,
]{article}
\usepackage{amsmath,amssymb}
\usepackage{setspace}
\usepackage{iftex}
\ifPDFTeX
  \usepackage[T1]{fontenc}
  \usepackage[utf8]{inputenc}
  \usepackage{textcomp} % provide euro and other symbols
\else % if luatex or xetex
  \usepackage{unicode-math} % this also loads fontspec
  \defaultfontfeatures{Scale=MatchLowercase}
  \defaultfontfeatures[\rmfamily]{Ligatures=TeX,Scale=1}
\fi
\usepackage{lmodern}
\ifPDFTeX\else
  % xetex/luatex font selection
  \setmainfont[]{Liberation Serif}
  \setsansfont[]{Liberation Sans}
  \setmonofont[]{DejaVu Sans Mono}
\fi
% Use upquote if available, for straight quotes in verbatim environments
\IfFileExists{upquote.sty}{\usepackage{upquote}}{}
\IfFileExists{microtype.sty}{% use microtype if available
  \usepackage[]{microtype}
  \UseMicrotypeSet[protrusion]{basicmath} % disable protrusion for tt fonts
}{}
\makeatletter
\@ifundefined{KOMAClassName}{% if non-KOMA class
  \IfFileExists{parskip.sty}{%
    \usepackage{parskip}
  }{% else
    \setlength{\parindent}{0pt}
    \setlength{\parskip}{6pt plus 2pt minus 1pt}}
}{% if KOMA class
  \KOMAoptions{parskip=half}}
\makeatother
\usepackage{xcolor}
\usepackage[margin=22mm]{geometry}
\setlength{\emergencystretch}{3em} % prevent overfull lines
\providecommand{\tightlist}{%
  \setlength{\itemsep}{0pt}\setlength{\parskip}{0pt}}
\setcounter{secnumdepth}{-\maxdimen} % remove section numbering
\ifLuaTeX
\usepackage[bidi=basic]{babel}
\else
\usepackage[bidi=default]{babel}
\fi
\babelprovide[main,import]{korean}
\ifPDFTeX
\else
\babelfont{rm}[]{Liberation Serif}
\fi
% get rid of language-specific shorthands (see #6817):
\let\LanguageShortHands\languageshorthands
\def\languageshorthands#1{}
\usepackage{xeCJK}
\setCJKmainfont{Noto Serif CJK KR}
\setCJKsansfont{Noto Sans CJK KR}
\setCJKmonofont{Noto Sans Mono CJK KR}
\xeCJKsetup{CJKspace=true}
\usepackage{xurl}
\usepackage{fvextra}
\usepackage{fancyhdr}
\AtBeginDocument{\DefineVerbatimEnvironment{verbatim}{Verbatim}{breaklines=true,breakanywhere=true,fontsize=\small}}
\setlength{\emergencystretch}{4em}
\setlength{\headheight}{15pt}
\pagestyle{fancy}\fancyhf{}
\fancyhead[L]{\small 중성자별 구조와 비방사 진동}
\fancyfoot[C]{\thepage}
\renewcommand{\headrulewidth}{0.3pt}
\setcounter{tocdepth}{2}
\ifLuaTeX
  \usepackage{selnolig}  % disable illegal ligatures
\fi
\usepackage{bookmark}
\IfFileExists{xurl.sty}{\usepackage{xurl}}{} % add URL line breaks if available
\urlstyle{same}
\hypersetup{
  pdftitle={A08 · 해석적 shooting: root와 pole},
  pdflang={ko-KR},
  colorlinks=true,
  linkcolor={Maroon},
  filecolor={Maroon},
  citecolor={Blue},
  urlcolor={Blue},
  pdfcreator={LaTeX via pandoc}}

\title{A08 · 해석적 shooting: root와 pole}
\usepackage{etoolbox}
\makeatletter
\providecommand{\subtitle}[1]{% add subtitle to \maketitle
  \apptocmd{\@title}{\par {\large #1 \par}}{}{}
}
\makeatother
\subtitle{과제 큐시트}
\author{}
\date{}

\begin{document}
\maketitle

\setstretch{1.12}
\textbf{선행:} A04　 \textbf{독립 작업:} 5--8시간 (계획치)

\subsection{부호가 바뀌는 곳마다 고유값이
있는가?}\label{uxbd80uxd638uxac00-uxbc14uxb00cuxb294-uxacf3uxb9c8uxb2e4-uxace0uxc720uxac12uxc774-uxc788uxb294uxac00}

표면값을 그 도함수로 나누면 분모가 0인 곳에서도 부호가 바뀔 수 있다. 그
지점을 고유값으로 부를 수 있는지 해석해를 보기 전에 판단하자.

먼저 한두 문장으로 답한 뒤 아래 설명을 읽자. 첫 예측이 틀렸거나 아직
모르겠다고 적는 것은 감점하지 않는다.

\subsection{필요한 관계와 읽을
위치}\label{uxd544uxc694uxd55c-uxad00uxacc4uxc640-uxc77duxc744-uxc704uxce58}

\(y^{\prime\prime}+k^2y=0\), \(y(0)=y(L)=0\), \(k_n=n\pi/L\).
\(\mathcal R=y(L)/(|A|L)\)는 진폭에 안전하지만 \(y/y^{\prime}\)는 pole을
가진다.

\begin{itemize}
\tightlist
\item
  \path|notes/chapters/ch04.md|
\item
  \path|notes/equations/shooting.md|
\item
  \path|learner/src/modes/ns_shooting_toy.f90|
\item
  \path|learner/src/modes/ns_scan.f90|
\end{itemize}

\subsection{계산으로
옮기기}\label{uxacc4uxc0b0uxc73cuxb85c-uxc62euxae30uxae30}

기본 경로는 제공 함수의 호출과 해석이다. 아래 \texttt{learner} 시험은
같은 물리 함수를 직접 작성한 경우에 실행한다. 제공 예제를 검사할 때에는
명시적으로 \texttt{calculation}을 선택한다.

\begin{enumerate}
\def\labelenumi{\arabic{enumi}.}
\tightlist
\item
  A≠0와 k=0 예외를 해석적으로 확인한다.
\item
  single shot→scan→bracket→refine→재적분 순서로 구현한다.
\item
  실패 trial을 지우지 말고 bracket barrier로 남긴다.
\item
  L과 scan 밀도를 바꾸고 누락된 root 가능성을 조사한다.
\end{enumerate}

\newpage

\subsection{실행해서 무엇을 볼
것인가}\label{uxc2e4uxd589uxd574uxc11c-uxbb34uxc5c7uxc744-uxbcfc-uxac83uxc778uxac00}

패키지 루트에서 실행한다. \texttt{-\/-plan}은 경로만 확인하며 계산
성공을 뜻하지 않는다.

\begin{verbatim}
python tools/run_assignment.py A08 --plan
python tools/run_assignment.py A08 --run \
  --workspace learner --output ../A08_result.json
\end{verbatim}

\subsection{결과를
설명하기}\label{uxacb0uxacfcuxb97c-uxc124uxba85uxd558uxae30}

판단의 전제는 잔차의 연속성과 비자명한 해다. 실패한 trial을 지운 뒤
양옆을 잇거나 영해를 모드로 세면 다른 문제를 풀게 된다.

\textbf{남길 증거.} 잔차 scan 및 pole 반례; 두 고유함수와 k\_n L 비교;
ODE 오차와 root 종료오차의 분리; 양방향 해석해의 determinant가
-sin(kL)/k임을 확인한 계산

\textbf{판단 기준.} 2점 예: 실패한 trial 양옆을 연결하지 않는 이유를
연속성과 함께 설명한다. 최초 예측의 정오보다, 결과를 본 뒤 이유를 고치고
다른 입력에도 적용하는지를 본다.

\subsection{막혔을 때}\label{uxb9c9uxd614uxc744-uxb54c}

\textbf{힌트 1.} 한 trial의 출력은 경계조건을 아직 강제하지 않은 초기값
해이다.

\textbf{힌트 2.} π/2 근처 y/y′의 부호변화와 올바른 잔차의 부호를 나란히
적는다.

\subsection{다음 만남에 다시
묻기}\label{uxb2e4uxc74c-uxb9ccuxb0a8uxc5d0-uxb2e4uxc2dc-uxbb3buxae30}

다음 회차에는 길이를 바꾸되 새로운 정답을 먼저 보지 않고 두 근을
예측한다. 마지막에 무차원 곱으로 해석해와 대조한다. 먼저 짧게 답한 뒤
본문과 자신의 기록을 열어 확인한다.

\textbf{범위.} 이 toy의 k는 역길이이고 중성자별의 omega가 아니다. node
규칙을 sharp interface에 무검토 전사하지 않는다.

\end{document}
